\section*{Verbs}

Verbs are constructed around a consonant root that denotes its aspect. This root is combined with a suffix to indicate tense. The tense suffixes are \textit{-e} for past,\textit{ -a} for present, and \textit{-o} for future. Optional prefixes can be added to denote mood and voice. This final conjugated auxiliary verb is placed in front of the infinitive. The verb that indicates the action is never conjugated, only the auxiliary verb changes form to show person, aspect, tense, mood, and voice. The lexical predicate is an unchanged noun root used verbally. Every auxiliary requires a person prefix in formal writing and complex sentences. The order is person + mood + voice + aspect root + tense.

\begin{table}[H]
    \centering
    \begin{tabular}{|c|c|}
        \hline
        \textbf{Tense} & \textbf{Auxiliary Verb Suffix} \\
        \hline
        Past & -e \\
        Present & -a \\
        Future & -o \\
        \hline
    \end{tabular}
    % \caption{Person Suffixes in Tovian}
    \label{tab:verb-tense}
\end{table}


For example, \textit{fethr} is the noun root for speech, used verbally as ``speak.'' The continuous aspect root \textit{sh-} takes present \textit{-a} and third-person \textit{a-}, forming \textit{asha}. The auxiliary precedes the unchanged predicate: \textit{i-la-yel asha fethr}, ``the person is speaking.''

\begin{table}[h!]
    \centering
    \subsubsection{Aspects with present tense examples}
    \begin{tabular}{c|c|c|c}
    \hline
         \textbf{Aspect}& \makecell{\textbf{Root}\\\textbf{Verb}}& \makecell{\textbf{Aux.}\\\textbf{root}} & \makecell{\textbf{Pres. tense}\\\textbf{Example}} \\
\hline
         Simple&\makecell{\textit{fethr} \\ speak}  & f- & \makecell{\textit{afa muril}\\``sees"} \\ \hline
         \red{Imperfective}&\makecell{\textit{thraf} \\ wonder}  & th- & \makecell{\textit{atha muril}\\ ``is seeing"} \\ \hline
         Perfect&\makecell{\textit{lhara} \\ know}  & lh- & \makecell{\textit{alha muril} \\ ``has seen"} \\ \hline
         Near&\makecell{\textit{tlun} \\ share}& tl- & \makecell{\textit{atla muril} \\ ``is about to see" }\\ \hline
         Immediate&\makecell{\textit{kesh} \\ burn}& k- & \makecell{\textit{aka muril} \\ ``is seeing right now" }\\ \hline
         Habitual&\makecell{\textit{mel} \\ live}  & m- & \makecell{\textit{ama muril} \\``sees regularly" }\\ \hline
         Progressive&\makecell{\textit{sethr} \\ run}& s- & \makecell{\textit{asa muril} \\``is seeing now" }\\ \hline
         Continuous&\makecell{\textit{shotl} \\ flow}& sh- & \makecell{\textit{asha muril} \\ ``is seeing" }\\ \hline
         Iterative&\makecell{\textit{nalh} \\ think}& n- & \makecell{\textit{ana muril} \\ ``sees repeatedly" }\\ \hline
         Inceptive&\makecell{\textit{yalor} \\ begin}& y- & \makecell{\textit{aya muril} \\ ``begins  seeing" }\\ \hline
         Cessative&\makecell{\textit{panor} \\ end}& p- & \makecell{\textit{apa muril} \\ ``stops seeing" }\\ \hline
         Remote&\makecell{\textit{huysh} \\ throw}  & h- & \makecell{\textit{aha muril} \\``sees (distant)" \footnotemark[1]} \\
     \end{tabular}
     \\
     \medskip
     \raggedright
     {\footnotesize \footnotemark[1] \raggedleft In past and future tenses, the remote aspect refers to actions that took place long ago or far in the future, but the meaning in present tense is more nuanced. Elves tend to use the rare present remote aspect as an insult, for example \textit{ta-e aha fethr}: ``they speak but it is irrelevant"}
    \label{tab:aspects}
\end{table}

\medskip


\subsection*{Mood and Voice}

% \subsection*{All verb forms}
\begin{tabular}{rl}
&\textbf{Mood}\\
\textbf{unmarked}& $\rightarrow$ indicative \\
\textbf{wi-}& $\rightarrow$ subjunctive \\
\textbf{se-}& $\rightarrow$ imperative \\
\textbf{ko-} & $\rightarrow$ conditional \\
\textbf{xo-}& $\rightarrow$ counterfactual \\
\textbf{hle-}& $\rightarrow$ optative \\
\textbf{tlo-}& $\rightarrow$ obligative \\
\textbf{ef-}& $\rightarrow$ necessative \\
\textbf{sho-}& $\rightarrow$ potential \\



&\textbf{Voices}\\
\textbf{unmarked}& $\rightarrow$ active \\
\textbf{te-}& $\rightarrow$ reflexive \\
\textbf{pa-}& $\rightarrow$ passive \\
\textbf{mo-}& $\rightarrow$ middle \\
\textbf{ke-}& $\rightarrow$ causative \\
\textbf{ra-}& $\rightarrow$ reciprocal \\

% &\textbf{Other (not combined)}\\
% \textbf{li-}& $\rightarrow$ gerund \\

\end{tabular}

\medskip

Mood and voice can be combined by adding the mood prefix before the voice prefix. For example, adding the obligative mood marker \textit{tlo-} and the reflexive voice marker \textit{te-} to the habitual aspect auxilary root stem \textit{m-} with the present tense conjugation \textit{-a}, with third-person \textit{a-}, the final auxiliary verb is \textit{atlotema}.

\ex
\begingl
\gla i-la-yel atlotema sor //
\glb NOM-ANIM-person 3-OBLIG-REFL-HAB.PRES wash //
\glft "She should habitually wash herself." //
\endgl
\xe



\subsection{Unused Combinations}
Tovian is very flexible when it comes to combining different tenses, aspects, moods, and voices. However, certain verb conjugations result in combinations that are unlikely or never used due to their impractical or nonsensical nature. For example, the Past Imperative is very rare because commands most often refer to actions in the present or future.



\section{Copulas}
Copulas are formed by combining verbs and nouns to describe states and qualities. For example the word "mel" (life/living) can be used as a copula: "i-la-yel afa mel" means "The person is alive," or equally "The person lives" This method applies to various attributes by changing the verb root and noun. For instance, "yel afe yeol" translates to "The person was wise."

 \section{Comparative Constructions}

    Tovian uses comparative constructions to express excessiveness, for example: "bigger than appropriate." The ablative case is often used for this comparison: "the tree is too big" could be expressed as \textit{i-la-gelar \CAabl-le-velar afa mot} ("The tree is big away-from-worthiness" or simply "The tree is bigger than what is worthy/right").

\section{Adjectives}

Tovian does not use a separate class of adjectives. Instead, all descriptive qualities are expressed through nouns. These attributive nouns follow the noun they modify in a fixed noun–noun structure. For example, the phrase \textit{yel yeol} translates literally as “person wisdom,” conveying the meaning “the wise person.” To convey "people wisdom" or "the wisdom of people," one would switch the order to: \textit{yeol yel}. This structure does not need inflection or agreement between adjectives and nouns, for example "the wise people" would simply be \textit{yele yeol}. In cases requiring emphasis or clarification, a copular construction can be used instead: \textit{i-la-yel afa yeol} (“The person is wise”). Many commonly used descriptive terms—such as colors, emotional states, or physical traits—exist as standalone nouns and follow the same syntactic pattern. This noun-only system reinforces the conceptual unity of properties and entities, treating attributes not as modifiers but as qualities with independent substance.


\section{Agreement}

Person is marked by a prefix on every auxiliary in formal writing and complex sentences. Stating a subject noun or pronoun does not remove this requirement. Lexical predicates take no person suffix.

\begin{center}
\begin{tabular}{c|c|l}
\textbf{Person} & \textbf{Prefix} & \textbf{Example} \\\hline
1st & i- & \textit{isa nalh} (I/we think) \\
2nd & o- & \textit{osa nalh} (you think) \\
3rd & a- & \textit{asa nalh} (he/she/it/they think) \\
\end{tabular}
\end{center}

Number is expressed on nouns, not by auxiliary agreement. A plural subject such as ``the truths'' therefore uses the same third-person \textit{a-} as a singular subject.

In short, unambiguous casual speech only, third-person \textit{a-} may be omitted: \textit{fa muril} can stand for \textit{afa muril}. Retain person marking in formal writing and complex sentences, including every auxiliary in a chain. First- and second-person prefixes remain required.

Bare \textit{lha} means ``hardness/strength'' and can modify a preceding noun. The formal perfect auxiliary is \textit{alha}; the passive perfect is \textit{apalha}:
\begin{itemize}
\item \textit{i-le-yeol-e lha apafa nalh}: ``The hard truths are pondered.''
\item \textit{i-le-yeol-e apalha nalh}: ``The truths have been pondered.''
\end{itemize}
The person prefix identifies the auxiliary while the lexical root stays unchanged.

\section{Auxiliary Nesting and Scope}

An auxiliary may take the following auxiliary phrase as its complement. Each earlier auxiliary applies to the entire phrase after it. Read a chain from the inside outward. Brackets here show grouping; they add no spoken or written Tovian morphemes.

Within one shared-subject chain, every auxiliary carries the same person prefix. Each layer may contribute its own mood and aspect. Within a single auxiliary, mood takes scope over aspect. Case and class prefixes remain on noun phrases in their existing order.

\subsection*{Scope contrasts}

\begin{itemize}
\item \textit{apa muril}: stops seeing.
\item \textit{ashopa muril}: can [stop seeing].
\item \textit{apa [ashofa muril]}: stops [being able to see].
\item \textit{ashoya muril}: can [begin seeing].
\item \textit{aya [ashofa muril]}: begins [being able to see].
\item \textit{atlopa for}: should [stop writing].
\item \textit{apa [atlofa for]}: ceases [to be expected to write].
\item \textit{ipa [ishofa muril]}: I/we cease [being able to see].
\end{itemize}

Agreement repeats because the grammatical subject is shared. The moods and aspects need not match.

\subsection*{Can no longer be ignored}

The intended scope is cessation over possibility over a passive predicate. With \texttt{[IGNORE]} as an English placeholder for a lexical root still to be chosen:

\begin{quote}
\texttt{apa [ashopafa [IGNORE]]}
\end{quote}

\begin{itemize}
\item \textit{a-p-a}: third person + cessative + present.
\item \textit{a-sho-pa-f-a}: third person + potential + passive + simple + present.
\item \texttt{[IGNORE]}: the unchanged lexical predicate.
\end{itemize}

The truths remain the grammatical subject throughout. The inner phrase means ``can be ignored''; the outer cessative ends that possibility. ``Stop being ignored'' describes the end of actual ignoring, while ``can stop being ignored'' expresses the possibility of stopping. The chain uses one cessative; the bound \textit{pa-} inside \textit{ashopafa} marks passive voice.

``No longer'' presupposes an earlier possibility without implying that the impossibility lasts forever. This construction adds neither a new particle nor an untensed auxiliary form.

\subsection*{Chains and embedded clauses}

``I cease [being able to see]'' has one shared subject. ``I know [you left yesterday]'' contains a clause with its own subject and time. Matching person is required within one shared-subject chain, not across an entire sentence. The temporal frame \textit{tu-le-ngothr: P; Q} specifies one construction with a finite event description P (see Temporal Clauses). General complement clauses such as ``I know [you left]'' still need their own rule.

\subsection*{Temporal dependence: open question}

The examples here use present tense throughout. Matching person does not require Tovian to copy tense endings across every layer. A general rule for inner temporal reference and mixed-tense chains remains to be defined. These examples establish person agreement and scope without settling temporal dependence.


%  Verb stems are added to nouns and always end in ``lore". Different tenses, aspects, moods, and voices can be applied by adding prefixes. The prefixes in the order they should be applied are as follows:\footnote{\noindent Complex verb forms are saved for formal communication.}

% \subsection{Simple Verb Forms}

% \begin{tabular}{rl}
% &\textbf{Tense}\\
% \textbf{} &$\rightarrow$$ present tense \\
% \textbf{vi-} &$\rightarrow$$ future tense \\
% \textbf{in-} &$\rightarrow$$ past tense \\

% &\textbf{Aspect}\\
% \textbf{de-}& $\rightarrow$ perfect \\
% \textbf{me-} & $\rightarrow$ participle \\
% \textbf{shi-}& $\rightarrow$ continuous \\

% &\textbf{Mood}\\
% \textbf{ko-} & $\rightarrow$ conditional \\
% \textbf{se-}& $\rightarrow$ imperative \\
% \textbf{wi-}& $\rightarrow$ subjunctive \\

% &\textbf{Voices}\\
% \textbf{te-}& $\rightarrow$ reflexive \\
% \textbf{pe-}& $\rightarrow$ passive \\

% &\textbf{Other (not combined)}\\
% \textbf{li-}& $\rightarrow$ gerund \\


% \end{tabular}
% % 55 combinations


% \subsection*{All verb forms}
% \begin{tabular}{rl}
% &\textbf{Tense}\\
% \textbf{} &$\rightarrow$$ present tense \\
% \textbf{vi-} &$\rightarrow$$ future tense \\
% \textbf{in-} &$\rightarrow$$ past tense \\
% &\textbf{Aspect}\\
% \textbf{de-}& $\rightarrow$ perfect \\
% \textbf{me-} & $\rightarrow$ participle \\
% \textbf{re-} & $\rightarrow$ habitual \\
% \textbf{ro-}& $\rightarrow$ iterative \\
% \textbf{ve-} & $\rightarrow$ immediate \\
% \textbf{shi-}& $\rightarrow$ continuous \\
% \textbf{wo-}& $\rightarrow$ remote \\
% \textbf{to-}& $\rightarrow$ near \\
% \textbf{yi-}& $\rightarrow$ inceptive \\
% \textbf{hlo-}& $\rightarrow$ cessative \\

% &\textbf{Mood}\\
% \textbf{ko-} & $\rightarrow$ conditional \\
% \textbf{se-}& $\rightarrow$ imperative \\
% \textbf{wi-}& $\rightarrow$ subjunctive \\
% \textbf{xo-}& $\rightarrow$ counterfactual \\

% &\textbf{Voices}\\
% \textbf{te-}& $\rightarrow$ reflexive \\
% \textbf{pa-}& $\rightarrow$ passive \\
% \textbf{mo-}& $\rightarrow$ middle \\
% \textbf{ke-}& $\rightarrow$ causative \\
% \textbf{ra-}& $\rightarrow$ reciprocal \\

% &\textbf{Other (not combined)}\\
% \textbf{li-}& $\rightarrow$ gerund \\


% \end{tabular}
% % 601 combinations

% \medskip

% \subsection*{Verb Conjugations}
% Using "speak" as an example: \textit{fethr}

% \subsubsection*{Tense}
% \begin{itemize}
%     \item \textbf{Present Tense:} \textit{} \\
%           \hanging Speak: \textit{lore-fethr}

%     \item \textbf{Future Tense:} \textit{vi-} \\
%           \hanging Will speak: \textit{vilore-fethr}

%     \item \textbf{Past Tense:} \textit{in-} \\
%           \hanging Spoke: \textit{inlore-fethr}
% \end{itemize}

% \subsubsection*{Aspect}
% \begin{itemize}
%     \item \textbf{Perfect:} \textit{de-} \\
%           \hanging Has spoken: \textit{delore-fethr}

%     \item \textbf{Continuous:} \textit{shi-} \\
%           \hanging Is speaking: \textit{shilore-fethr}

%     \item \textbf{Habitual:} \textit{re-} \\
%           \hanging Speaks regularly: \textit{relore-fethr}

%     \item \textbf{Immediate:} \textit{ve-} \\
%           \hanging Is about to speak: \textit{velore-fethr}

%     \item \textbf{Participle:} \textit{ma-} \\
%           \hanging Speaking (as a participle): \textit{malore-fethr}
% \end{itemize}

% \subsubsection*{Mood}
% \begin{itemize}
%     \item \textbf{Conditional:} \textit{ko-} \\
%           \hanging Would speak: \textit{kolore-fethr}

%     \item \textbf{Imperative:} \textit{se-} \\
%           \hanging Speak!: \textit{selore-fethr}

%     \item \textbf{Subjunctive:} \textit{wi-} \\
%           \hanging May speak: \textit{wilore-fethr}

%     \item \textbf{Counterfactual:} \textit{xo-} \\
%           \hanging Would have spoken: \textit{xolore-fethr}
% \end{itemize}

% \subsubsection*{Voice}
% \begin{itemize}
%     \item \textbf{Reflexive:} \textit{te-} \\
%           \hanging Speaks to oneself: \textit{telore-fethr}

%     \item \textbf{Passive:} \textit{pa-} \\
%           \hanging Is spoken: \textit{palore-fethr}

%     \item \textbf{Middle:} \textit{mo-} \\
%           \hanging Speaks for oneself: \textit{molore-fethr}

%     \item \textbf{Causative:} \textit{ke-} \\
%           \hanging Makes someone speak: \textit{kelore-fethr}

%     \item \textbf{Reciprocal:} \textit{ra-} \\
%           \hanging They speak to each other: \textit{ralore-fethr}
% \end{itemize}

% \subsubsection*{Other}
% \begin{itemize}
%     \item \textbf{Gerund:} \textit{li-} \\
%           \hanging Speaking (as a noun): \textit{lilore-fethr}
% \end{itemize}

% \subsubsection*{Combined Forms}
% \begin{itemize}
%     \item \textbf{Past Perfect:} \textit{de- + in-} \\
%           \hanging Had spoken: \textit{deinlore-fethr}

%     \item \textbf{Future Continuous:} \textit{shi- + vi-} \\
%           \hanging Will be speaking: \textit{shivilore-fethr}

%     \item \textbf{Conditional Habitual:} \textit{re- + ko-} \\
%           \hanging Would speak regularly: \textit{rekolore-fethr}

%     \item \textbf{Immediate Imperative:} \textit{ve- + se-} \\
%           \hanging Speak immediately!: \textit{veselore-fethr}

%     \item \textbf{Subjunctive Participle:} \textit{wi- + ma-} \\
%           \hanging Speaking in a hypothetical or wished-for sense: \textit{wimalore-fethr}
% \end{itemize}


% \subsection*{To be and to be like}
% All verbs with the exception of two follow the above structure. To express that the subject of the sentence is equivalent to the object, use ``'as" as a postfix on the object. To express similarity, use ``'asl" instead. For example:

% \begin{itemize}
%   \item \textbf{To be:} \\
%         \textit{Tel \CAacc-\CLin-helor'as.} (The lake is glass.) \\
%         - \textit{Tel} (subject: the lake) \\
%         - \textit{\CAacc-\CLin-helor'as} (object: glass with ``as" indicating "is")

%  \item \textbf{To be like:} \\
%         \textit{Tel \CAacc-\CLin-helor'asl.} (The lake is like glass.) \\
%         - \textit{Tel} (subject: the lake) \\
%         - \textit{helor'as} (object: glass with ``asl'" indicating "is like")

% \end{itemize}

